% ------------------------------------------------------------------------
%  FMM implementation: whole-pipeline overview, group headlines only.
%  Include with % ------------------------------------------------------------------------
%  FMM implementation: whole-pipeline overview, group headlines only.
%  Include with % ------------------------------------------------------------------------
%  FMM implementation: whole-pipeline overview, group headlines only.
%  Include with % ------------------------------------------------------------------------
%  FMM implementation: whole-pipeline overview, group headlines only.
%  Include with \input{figures/fmm_overview}
%  Each group is broken out in its own detailed figure (fmm_g1 .. fmm_g7).
% ------------------------------------------------------------------------
\input{figures/fmm_style}

\begin{figure}[p]
\centering
\begin{tikzpicture}[
    >={Latex[length=1.9mm]},
    b/.style ={fmmgpu, text width=4.4cm, align=center, inner sep=1.4mm,
               font=\small},
    bc/.style={fmmcpu, text width=4.4cm, align=center, inner sep=1.4mm,
               font=\small},
    w/.style ={b,  text width=9.7cm},
    wc/.style={bc, text width=9.7cm},
]

% ---------- legend ----------
\node[anchor=north, align=center] at (4.99,0.5) {{\scriptsize
   \tikz[baseline=-0.5ex]\node[draw, line width=0.6pt, fill=white,
        minimum width=4mm, minimum height=2.6mm, inner sep=0pt]{};\;GPU kernel
   \qquad
   \tikz[baseline=-0.5ex]\node[draw=black!55, line width=0.5pt,
        dash pattern=on 1.4pt off 1.2pt, fill=black!7,
        minimum width=4mm, minimum height=2.6mm, inner sep=0pt]{};\;host (CPU) step}};

% =============== 1  source tree construction ===============
\node[wc, anchor=north west] (a1) at (0,-1.0) {Setup};
\node[w,  anchor=north west] (a2) at ([yshift=-1.8mm]a1.south west) {LBVH build};
\node[w,  anchor=north west] (a3) at ([yshift=-1.8mm]a2.south west) {Box nodes and box leaves};
\node[wc, anchor=north west] (a4) at ([yshift=-1.8mm]a3.south west) {Moment storage};
\node[fmmgrp, fit=(a1)(a2)(a3)(a4)] (G1) {};

% =============== 2  upward pass ===============
% Every group is anchored back to x = 0, never to the previous PANEL edge --
% otherwise each group drifts 2.4 mm (one fmmgrp inner sep) further left.
\coordinate (p2) at (0,0 |- G1.south);
\node[b, anchor=north west] (b1) at ([yshift=-8mm]p2) {P2M};
\node[b, anchor=north west] (b2) at ([xshift=6mm]b1.north east) {M2M};
\node[fmmgrp, fit=(b1)(b2)] (G2) {};

% =============== 3  target tree construction ===============
\coordinate (p3) at (0,0 |- G2.south);
\node[w,  anchor=north west] (c1) at ([yshift=-23mm]p3) {LBVH build, box nodes and box leaves};
\node[w,  anchor=north west] (c2) at ([yshift=-1.8mm]c1.south west) {L2L depth worklist};
\node[wc, anchor=north west] (c3) at ([yshift=-1.8mm]c2.south west) {Local-expansion storage};
\node[fmmgrp, fit=(c1)(c2)(c3)] (G3) {};

% =============== 4  dual-tree traversal ===============
\coordinate (p4) at (0,0 |- G3.south);
\node[bc, anchor=north west] (d1) at ([yshift=-8mm]p4) {Host driver};
\node[b,  anchor=north west] (d2) at ([xshift=6mm]d1.north east) {\texttt{\small fmm\_pair\_kernel}};
\node[fmmgrp, fit=(d1)(d2)] (G4) {};

% =============== 5 / 6  the two interaction lists ===============
% M2L and L2L share one host function (fmm_translations<P>), so they share a panel.
\coordinate (p5) at (0,0 |- G4.south);
\node[b, anchor=north west] (e1) at ([yshift=-11mm]p5) {Near-list CSR};
\node[b, anchor=north west] (f1) at ([xshift=6mm]e1.north east) {M2L accumulate};
\node[b, anchor=north west] (f2) at ([yshift=-1.8mm]f1.south west) {L2L level sweep};
\node[fmmgrp, fit=(e1)] (G5) {};
\node[fmmgrp, fit=(f1)(f2)] (G6) {};

% =============== 7  fused evaluation ===============
\coordinate (p7) at (0,0 |- G6.south);
\node[w,  anchor=north west] (h1) at ([yshift=-11mm]p7) {Fused L2P $+$ P2P leaf kernel};
\node[wc, anchor=north west] (h2) at ([yshift=-1.8mm]h1.south west) {Host reduction};
\node[fmmgrp, fit=(h1)(h2)] (G7) {};

% ---------- group labels on the panel borders ----------
\node[fmmlbl] at ([xshift=4mm]G1.north west) {\bfseries 1\ \ Source tree construction};
\node[fmmlbl] at ([xshift=4mm]G2.north west) {\bfseries 2\ \ Upward pass};
\node[fmmlbl] at ([xshift=4mm]G3.north west) {\bfseries 3\ \ Target tree construction};
\node[fmmlbl] at ([xshift=4mm]G4.north west) {\bfseries 4\ \ Dual-tree traversal};
\node[fmmlbl] at ([xshift=4mm]G5.north west) {\bfseries 5\ \ P2P CSR};
\node[fmmlbl] at ([xshift=4mm]G6.north west) {\bfseries 6\ \ Downward pass};
\node[fmmlbl] at ([xshift=4mm]G7.north west) {\bfseries 7\ \ Fused evaluation};

% ---------- flow ----------
\foreach \x/\y in {a1/a2, a2/a3, a3/a4, c1/c2, c2/c3, f1/f2, h1/h2}
  \draw[fmmflow] (\x.south) -- (\y.north);
\draw[fmmflow] (G1.south) -- (G2.north);
\draw[fmmflow] (G2.south) -- (G3.north);
\draw[fmmflow] (G3.south) -- (G4.north);

% Connectors run down the right of each panel: the group labels sit on the
% top-left of the borders, so a centered arrow would strike through them.
\coordinate (x5) at ([xshift=-8mm]G5.north east);
\coordinate (x6) at ([xshift=-8mm]G6.north east);

% the single traversal pass forks into the two interaction lists
\draw[fmmflow] (x5 |- G4.south) -- node[fmmedge, pos=0.55] {P2P pair list} (x5);
\draw[fmmflow] (x6 |- G4.south) -- node[fmmedge, pos=0.55] {M2L pair list} (x6);

% and both rejoin in the fused leaf kernel
\draw[fmmflow] (x5 |- G5.south) -- node[fmmedge, pos=0.5] {near field} (x5 |- G7.north);
\draw[fmmflow] (x6 |- G6.south) -- node[fmmedge, pos=0.5] {far field} (x6 |- G7.north);

% ---------- lifecycle band (drawn last so it masks the arrow behind it) ------
\draw[black!45, line width=0.5pt, dash pattern=on 1.6pt off 1.6pt]
     ([xshift=-6mm,yshift=9.5mm]G3.north west) -- ([xshift=6mm,yshift=9.5mm]G3.north east);
\node[fmmband, anchor=south west] at ([xshift=-4mm,yshift=10.1mm]G3.north west)
     {once per solve};
\node[fmmband, anchor=north west] at ([xshift=-4mm,yshift=8.9mm]G3.north west)
     {once \emph{per energy term}: $E_{\text{pol}}$, then $E_{\text{ion}}$};

\end{tikzpicture}
\caption[FMM implementation: pipeline overview]{%
    \textbf{Overview of the FMM implementation.} Solid boxes are GPU kernels,
    dashed boxes host steps. Each group is broken out at the level of individual
    kernels in Figures~\ref{fig:fmm-g1}--\ref{fig:fmm-g7}.}
\label{fig:fmm-overview}
\end{figure}

%  Each group is broken out in its own detailed figure (fmm_g1 .. fmm_g7).
% ------------------------------------------------------------------------
% ------------------------------------------------------------------------
%  Shared styling for the FMM implementation figures:
%    fmm_overview.tex   -- the whole pipeline, group headlines only
%    fmm_g1..g8.tex     -- one detailed figure per group
%  Idempotent: every figure \input{figures/fmm_style} before its tikzpicture.
%  Needs tikz with the libraries positioning, fit and arrows.meta.
% ------------------------------------------------------------------------

% Entry macros. The mapping runs on from the kernel name in the same paragraph
% rather than on its own line -- a forced break per entry doubles the height.
\providecommand{\fmmH}[1]{{\footnotesize\bfseries #1}\par\vspace{0.5ex}}
\providecommand{\fmmK}[2]{{\ttfamily\scriptsize #1}\,{\tiny\itshape\color{black!62}#2}\par\vspace{0.35ex}}
\providecommand{\fmmN}[1]{{\tiny\color{black!62}#1}\par\vspace{0.35ex}}

% Legend, shared by the overview and by any detail figure that wants it.
\providecommand{\fmmlegend}{{\scriptsize
   \tikz[baseline=-0.5ex]\node[draw, line width=0.6pt, fill=white,
        minimum width=4mm, minimum height=2.6mm, inner sep=0pt]{};\;GPU kernel
   \quad
   \tikz[baseline=-0.5ex]\node[draw=black!55, line width=0.5pt,
        dash pattern=on 1.4pt off 1.2pt, fill=black!7,
        minimum width=4mm, minimum height=2.6mm, inner sep=0pt]{};\;host (CPU) step
   \quad {\itshape italics} $=$ parallel mapping
   \quad \texttt{cub::} $=$ CUB device-wide primitive (GPU)}}

\tikzset{
  fmmgpu/.style  ={draw, line width=0.6pt, fill=white, align=left, inner sep=1.5mm},
  fmmcpu/.style  ={fmmgpu, draw=black!55, dash pattern=on 1.4pt off 1.2pt,
                   fill=black!7},
  fmmnote/.style ={fmmgpu, draw=black!45, fill=black!4, rounded corners=2pt},
  fmmtag/.style  ={font=\tiny\sffamily, text=black!55, fill=white,
                   inner sep=0.5mm, anchor=north east},
  fmmflow/.style ={->, line width=0.7pt, black!70},
  fmmgrp/.style  ={draw=black!45, line width=0.9pt, rounded corners=2.5pt,
                   inner sep=2.4mm},
  fmmlbl/.style  ={font=\small, fill=white, inner sep=1mm, anchor=west},
  fmmedge/.style ={font=\tiny\sffamily, text=black!60, inner sep=0.7mm,
                   fill=white},
  fmmband/.style ={font=\scriptsize\itshape, text=black!60, inner sep=1mm,
                   fill=white},
}


\begin{figure}[p]
\centering
\begin{tikzpicture}[
    >={Latex[length=1.9mm]},
    b/.style ={fmmgpu, text width=4.4cm, align=center, inner sep=1.4mm,
               font=\small},
    bc/.style={fmmcpu, text width=4.4cm, align=center, inner sep=1.4mm,
               font=\small},
    w/.style ={b,  text width=9.7cm},
    wc/.style={bc, text width=9.7cm},
]

% ---------- legend ----------
\node[anchor=north, align=center] at (4.99,0.5) {{\scriptsize
   \tikz[baseline=-0.5ex]\node[draw, line width=0.6pt, fill=white,
        minimum width=4mm, minimum height=2.6mm, inner sep=0pt]{};\;GPU kernel
   \qquad
   \tikz[baseline=-0.5ex]\node[draw=black!55, line width=0.5pt,
        dash pattern=on 1.4pt off 1.2pt, fill=black!7,
        minimum width=4mm, minimum height=2.6mm, inner sep=0pt]{};\;host (CPU) step}};

% =============== 1  source tree construction ===============
\node[wc, anchor=north west] (a1) at (0,-1.0) {Setup};
\node[w,  anchor=north west] (a2) at ([yshift=-1.8mm]a1.south west) {LBVH build};
\node[w,  anchor=north west] (a3) at ([yshift=-1.8mm]a2.south west) {Box nodes and box leaves};
\node[wc, anchor=north west] (a4) at ([yshift=-1.8mm]a3.south west) {Moment storage};
\node[fmmgrp, fit=(a1)(a2)(a3)(a4)] (G1) {};

% =============== 2  upward pass ===============
% Every group is anchored back to x = 0, never to the previous PANEL edge --
% otherwise each group drifts 2.4 mm (one fmmgrp inner sep) further left.
\coordinate (p2) at (0,0 |- G1.south);
\node[b, anchor=north west] (b1) at ([yshift=-8mm]p2) {P2M};
\node[b, anchor=north west] (b2) at ([xshift=6mm]b1.north east) {M2M};
\node[fmmgrp, fit=(b1)(b2)] (G2) {};

% =============== 3  target tree construction ===============
\coordinate (p3) at (0,0 |- G2.south);
\node[w,  anchor=north west] (c1) at ([yshift=-23mm]p3) {LBVH build, box nodes and box leaves};
\node[w,  anchor=north west] (c2) at ([yshift=-1.8mm]c1.south west) {L2L depth worklist};
\node[wc, anchor=north west] (c3) at ([yshift=-1.8mm]c2.south west) {Local-expansion storage};
\node[fmmgrp, fit=(c1)(c2)(c3)] (G3) {};

% =============== 4  dual-tree traversal ===============
\coordinate (p4) at (0,0 |- G3.south);
\node[bc, anchor=north west] (d1) at ([yshift=-8mm]p4) {Host driver};
\node[b,  anchor=north west] (d2) at ([xshift=6mm]d1.north east) {\texttt{\small fmm\_pair\_kernel}};
\node[fmmgrp, fit=(d1)(d2)] (G4) {};

% =============== 5 / 6  the two interaction lists ===============
% M2L and L2L share one host function (fmm_translations<P>), so they share a panel.
\coordinate (p5) at (0,0 |- G4.south);
\node[b, anchor=north west] (e1) at ([yshift=-11mm]p5) {Near-list CSR};
\node[b, anchor=north west] (f1) at ([xshift=6mm]e1.north east) {M2L accumulate};
\node[b, anchor=north west] (f2) at ([yshift=-1.8mm]f1.south west) {L2L level sweep};
\node[fmmgrp, fit=(e1)] (G5) {};
\node[fmmgrp, fit=(f1)(f2)] (G6) {};

% =============== 7  fused evaluation ===============
\coordinate (p7) at (0,0 |- G6.south);
\node[w,  anchor=north west] (h1) at ([yshift=-11mm]p7) {Fused L2P $+$ P2P leaf kernel};
\node[wc, anchor=north west] (h2) at ([yshift=-1.8mm]h1.south west) {Host reduction};
\node[fmmgrp, fit=(h1)(h2)] (G7) {};

% ---------- group labels on the panel borders ----------
\node[fmmlbl] at ([xshift=4mm]G1.north west) {\bfseries 1\ \ Source tree construction};
\node[fmmlbl] at ([xshift=4mm]G2.north west) {\bfseries 2\ \ Upward pass};
\node[fmmlbl] at ([xshift=4mm]G3.north west) {\bfseries 3\ \ Target tree construction};
\node[fmmlbl] at ([xshift=4mm]G4.north west) {\bfseries 4\ \ Dual-tree traversal};
\node[fmmlbl] at ([xshift=4mm]G5.north west) {\bfseries 5\ \ P2P CSR};
\node[fmmlbl] at ([xshift=4mm]G6.north west) {\bfseries 6\ \ Downward pass};
\node[fmmlbl] at ([xshift=4mm]G7.north west) {\bfseries 7\ \ Fused evaluation};

% ---------- flow ----------
\foreach \x/\y in {a1/a2, a2/a3, a3/a4, c1/c2, c2/c3, f1/f2, h1/h2}
  \draw[fmmflow] (\x.south) -- (\y.north);
\draw[fmmflow] (G1.south) -- (G2.north);
\draw[fmmflow] (G2.south) -- (G3.north);
\draw[fmmflow] (G3.south) -- (G4.north);

% Connectors run down the right of each panel: the group labels sit on the
% top-left of the borders, so a centered arrow would strike through them.
\coordinate (x5) at ([xshift=-8mm]G5.north east);
\coordinate (x6) at ([xshift=-8mm]G6.north east);

% the single traversal pass forks into the two interaction lists
\draw[fmmflow] (x5 |- G4.south) -- node[fmmedge, pos=0.55] {P2P pair list} (x5);
\draw[fmmflow] (x6 |- G4.south) -- node[fmmedge, pos=0.55] {M2L pair list} (x6);

% and both rejoin in the fused leaf kernel
\draw[fmmflow] (x5 |- G5.south) -- node[fmmedge, pos=0.5] {near field} (x5 |- G7.north);
\draw[fmmflow] (x6 |- G6.south) -- node[fmmedge, pos=0.5] {far field} (x6 |- G7.north);

% ---------- lifecycle band (drawn last so it masks the arrow behind it) ------
\draw[black!45, line width=0.5pt, dash pattern=on 1.6pt off 1.6pt]
     ([xshift=-6mm,yshift=9.5mm]G3.north west) -- ([xshift=6mm,yshift=9.5mm]G3.north east);
\node[fmmband, anchor=south west] at ([xshift=-4mm,yshift=10.1mm]G3.north west)
     {once per solve};
\node[fmmband, anchor=north west] at ([xshift=-4mm,yshift=8.9mm]G3.north west)
     {once \emph{per energy term}: $E_{\text{pol}}$, then $E_{\text{ion}}$};

\end{tikzpicture}
\caption[FMM implementation: pipeline overview]{%
    \textbf{Overview of the FMM implementation.} Solid boxes are GPU kernels,
    dashed boxes host steps. Each group is broken out at the level of individual
    kernels in Figures~\ref{fig:fmm-g1}--\ref{fig:fmm-g7}.}
\label{fig:fmm-overview}
\end{figure}

%  Each group is broken out in its own detailed figure (fmm_g1 .. fmm_g7).
% ------------------------------------------------------------------------
% ------------------------------------------------------------------------
%  Shared styling for the FMM implementation figures:
%    fmm_overview.tex   -- the whole pipeline, group headlines only
%    fmm_g1..g8.tex     -- one detailed figure per group
%  Idempotent: every figure % ------------------------------------------------------------------------
%  Shared styling for the FMM implementation figures:
%    fmm_overview.tex   -- the whole pipeline, group headlines only
%    fmm_g1..g8.tex     -- one detailed figure per group
%  Idempotent: every figure \input{figures/fmm_style} before its tikzpicture.
%  Needs tikz with the libraries positioning, fit and arrows.meta.
% ------------------------------------------------------------------------

% Entry macros. The mapping runs on from the kernel name in the same paragraph
% rather than on its own line -- a forced break per entry doubles the height.
\providecommand{\fmmH}[1]{{\footnotesize\bfseries #1}\par\vspace{0.5ex}}
\providecommand{\fmmK}[2]{{\ttfamily\scriptsize #1}\,{\tiny\itshape\color{black!62}#2}\par\vspace{0.35ex}}
\providecommand{\fmmN}[1]{{\tiny\color{black!62}#1}\par\vspace{0.35ex}}

% Legend, shared by the overview and by any detail figure that wants it.
\providecommand{\fmmlegend}{{\scriptsize
   \tikz[baseline=-0.5ex]\node[draw, line width=0.6pt, fill=white,
        minimum width=4mm, minimum height=2.6mm, inner sep=0pt]{};\;GPU kernel
   \quad
   \tikz[baseline=-0.5ex]\node[draw=black!55, line width=0.5pt,
        dash pattern=on 1.4pt off 1.2pt, fill=black!7,
        minimum width=4mm, minimum height=2.6mm, inner sep=0pt]{};\;host (CPU) step
   \quad {\itshape italics} $=$ parallel mapping
   \quad \texttt{cub::} $=$ CUB device-wide primitive (GPU)}}

\tikzset{
  fmmgpu/.style  ={draw, line width=0.6pt, fill=white, align=left, inner sep=1.5mm},
  fmmcpu/.style  ={fmmgpu, draw=black!55, dash pattern=on 1.4pt off 1.2pt,
                   fill=black!7},
  fmmnote/.style ={fmmgpu, draw=black!45, fill=black!4, rounded corners=2pt},
  fmmtag/.style  ={font=\tiny\sffamily, text=black!55, fill=white,
                   inner sep=0.5mm, anchor=north east},
  fmmflow/.style ={->, line width=0.7pt, black!70},
  fmmgrp/.style  ={draw=black!45, line width=0.9pt, rounded corners=2.5pt,
                   inner sep=2.4mm},
  fmmlbl/.style  ={font=\small, fill=white, inner sep=1mm, anchor=west},
  fmmedge/.style ={font=\tiny\sffamily, text=black!60, inner sep=0.7mm,
                   fill=white},
  fmmband/.style ={font=\scriptsize\itshape, text=black!60, inner sep=1mm,
                   fill=white},
}
 before its tikzpicture.
%  Needs tikz with the libraries positioning, fit and arrows.meta.
% ------------------------------------------------------------------------

% Entry macros. The mapping runs on from the kernel name in the same paragraph
% rather than on its own line -- a forced break per entry doubles the height.
\providecommand{\fmmH}[1]{{\footnotesize\bfseries #1}\par\vspace{0.5ex}}
\providecommand{\fmmK}[2]{{\ttfamily\scriptsize #1}\,{\tiny\itshape\color{black!62}#2}\par\vspace{0.35ex}}
\providecommand{\fmmN}[1]{{\tiny\color{black!62}#1}\par\vspace{0.35ex}}

% Legend, shared by the overview and by any detail figure that wants it.
\providecommand{\fmmlegend}{{\scriptsize
   \tikz[baseline=-0.5ex]\node[draw, line width=0.6pt, fill=white,
        minimum width=4mm, minimum height=2.6mm, inner sep=0pt]{};\;GPU kernel
   \quad
   \tikz[baseline=-0.5ex]\node[draw=black!55, line width=0.5pt,
        dash pattern=on 1.4pt off 1.2pt, fill=black!7,
        minimum width=4mm, minimum height=2.6mm, inner sep=0pt]{};\;host (CPU) step
   \quad {\itshape italics} $=$ parallel mapping
   \quad \texttt{cub::} $=$ CUB device-wide primitive (GPU)}}

\tikzset{
  fmmgpu/.style  ={draw, line width=0.6pt, fill=white, align=left, inner sep=1.5mm},
  fmmcpu/.style  ={fmmgpu, draw=black!55, dash pattern=on 1.4pt off 1.2pt,
                   fill=black!7},
  fmmnote/.style ={fmmgpu, draw=black!45, fill=black!4, rounded corners=2pt},
  fmmtag/.style  ={font=\tiny\sffamily, text=black!55, fill=white,
                   inner sep=0.5mm, anchor=north east},
  fmmflow/.style ={->, line width=0.7pt, black!70},
  fmmgrp/.style  ={draw=black!45, line width=0.9pt, rounded corners=2.5pt,
                   inner sep=2.4mm},
  fmmlbl/.style  ={font=\small, fill=white, inner sep=1mm, anchor=west},
  fmmedge/.style ={font=\tiny\sffamily, text=black!60, inner sep=0.7mm,
                   fill=white},
  fmmband/.style ={font=\scriptsize\itshape, text=black!60, inner sep=1mm,
                   fill=white},
}


\begin{figure}[p]
\centering
\begin{tikzpicture}[
    >={Latex[length=1.9mm]},
    b/.style ={fmmgpu, text width=4.4cm, align=center, inner sep=1.4mm,
               font=\small},
    bc/.style={fmmcpu, text width=4.4cm, align=center, inner sep=1.4mm,
               font=\small},
    w/.style ={b,  text width=9.7cm},
    wc/.style={bc, text width=9.7cm},
]

% ---------- legend ----------
\node[anchor=north, align=center] at (4.99,0.5) {{\scriptsize
   \tikz[baseline=-0.5ex]\node[draw, line width=0.6pt, fill=white,
        minimum width=4mm, minimum height=2.6mm, inner sep=0pt]{};\;GPU kernel
   \qquad
   \tikz[baseline=-0.5ex]\node[draw=black!55, line width=0.5pt,
        dash pattern=on 1.4pt off 1.2pt, fill=black!7,
        minimum width=4mm, minimum height=2.6mm, inner sep=0pt]{};\;host (CPU) step}};

% =============== 1  source tree construction ===============
\node[wc, anchor=north west] (a1) at (0,-1.0) {Setup};
\node[w,  anchor=north west] (a2) at ([yshift=-1.8mm]a1.south west) {LBVH build};
\node[w,  anchor=north west] (a3) at ([yshift=-1.8mm]a2.south west) {Box nodes and box leaves};
\node[wc, anchor=north west] (a4) at ([yshift=-1.8mm]a3.south west) {Moment storage};
\node[fmmgrp, fit=(a1)(a2)(a3)(a4)] (G1) {};

% =============== 2  upward pass ===============
% Every group is anchored back to x = 0, never to the previous PANEL edge --
% otherwise each group drifts 2.4 mm (one fmmgrp inner sep) further left.
\coordinate (p2) at (0,0 |- G1.south);
\node[b, anchor=north west] (b1) at ([yshift=-8mm]p2) {P2M};
\node[b, anchor=north west] (b2) at ([xshift=6mm]b1.north east) {M2M};
\node[fmmgrp, fit=(b1)(b2)] (G2) {};

% =============== 3  target tree construction ===============
\coordinate (p3) at (0,0 |- G2.south);
\node[w,  anchor=north west] (c1) at ([yshift=-23mm]p3) {LBVH build, box nodes and box leaves};
\node[w,  anchor=north west] (c2) at ([yshift=-1.8mm]c1.south west) {L2L depth worklist};
\node[wc, anchor=north west] (c3) at ([yshift=-1.8mm]c2.south west) {Local-expansion storage};
\node[fmmgrp, fit=(c1)(c2)(c3)] (G3) {};

% =============== 4  dual-tree traversal ===============
\coordinate (p4) at (0,0 |- G3.south);
\node[bc, anchor=north west] (d1) at ([yshift=-8mm]p4) {Host driver};
\node[b,  anchor=north west] (d2) at ([xshift=6mm]d1.north east) {\texttt{\small fmm\_pair\_kernel}};
\node[fmmgrp, fit=(d1)(d2)] (G4) {};

% =============== 5 / 6  the two interaction lists ===============
% M2L and L2L share one host function (fmm_translations<P>), so they share a panel.
\coordinate (p5) at (0,0 |- G4.south);
\node[b, anchor=north west] (e1) at ([yshift=-11mm]p5) {Near-list CSR};
\node[b, anchor=north west] (f1) at ([xshift=6mm]e1.north east) {M2L accumulate};
\node[b, anchor=north west] (f2) at ([yshift=-1.8mm]f1.south west) {L2L level sweep};
\node[fmmgrp, fit=(e1)] (G5) {};
\node[fmmgrp, fit=(f1)(f2)] (G6) {};

% =============== 7  fused evaluation ===============
\coordinate (p7) at (0,0 |- G6.south);
\node[w,  anchor=north west] (h1) at ([yshift=-11mm]p7) {Fused L2P $+$ P2P leaf kernel};
\node[wc, anchor=north west] (h2) at ([yshift=-1.8mm]h1.south west) {Host reduction};
\node[fmmgrp, fit=(h1)(h2)] (G7) {};

% ---------- group labels on the panel borders ----------
\node[fmmlbl] at ([xshift=4mm]G1.north west) {\bfseries 1\ \ Source tree construction};
\node[fmmlbl] at ([xshift=4mm]G2.north west) {\bfseries 2\ \ Upward pass};
\node[fmmlbl] at ([xshift=4mm]G3.north west) {\bfseries 3\ \ Target tree construction};
\node[fmmlbl] at ([xshift=4mm]G4.north west) {\bfseries 4\ \ Dual-tree traversal};
\node[fmmlbl] at ([xshift=4mm]G5.north west) {\bfseries 5\ \ P2P CSR};
\node[fmmlbl] at ([xshift=4mm]G6.north west) {\bfseries 6\ \ Downward pass};
\node[fmmlbl] at ([xshift=4mm]G7.north west) {\bfseries 7\ \ Fused evaluation};

% ---------- flow ----------
\foreach \x/\y in {a1/a2, a2/a3, a3/a4, c1/c2, c2/c3, f1/f2, h1/h2}
  \draw[fmmflow] (\x.south) -- (\y.north);
\draw[fmmflow] (G1.south) -- (G2.north);
\draw[fmmflow] (G2.south) -- (G3.north);
\draw[fmmflow] (G3.south) -- (G4.north);

% Connectors run down the right of each panel: the group labels sit on the
% top-left of the borders, so a centered arrow would strike through them.
\coordinate (x5) at ([xshift=-8mm]G5.north east);
\coordinate (x6) at ([xshift=-8mm]G6.north east);

% the single traversal pass forks into the two interaction lists
\draw[fmmflow] (x5 |- G4.south) -- node[fmmedge, pos=0.55] {P2P pair list} (x5);
\draw[fmmflow] (x6 |- G4.south) -- node[fmmedge, pos=0.55] {M2L pair list} (x6);

% and both rejoin in the fused leaf kernel
\draw[fmmflow] (x5 |- G5.south) -- node[fmmedge, pos=0.5] {near field} (x5 |- G7.north);
\draw[fmmflow] (x6 |- G6.south) -- node[fmmedge, pos=0.5] {far field} (x6 |- G7.north);

% ---------- lifecycle band (drawn last so it masks the arrow behind it) ------
\draw[black!45, line width=0.5pt, dash pattern=on 1.6pt off 1.6pt]
     ([xshift=-6mm,yshift=9.5mm]G3.north west) -- ([xshift=6mm,yshift=9.5mm]G3.north east);
\node[fmmband, anchor=south west] at ([xshift=-4mm,yshift=10.1mm]G3.north west)
     {once per solve};
\node[fmmband, anchor=north west] at ([xshift=-4mm,yshift=8.9mm]G3.north west)
     {once \emph{per energy term}: $E_{\text{pol}}$, then $E_{\text{ion}}$};

\end{tikzpicture}
\caption[FMM implementation: pipeline overview]{%
    \textbf{Overview of the FMM implementation.} Solid boxes are GPU kernels,
    dashed boxes host steps. Each group is broken out at the level of individual
    kernels in Figures~\ref{fig:fmm-g1}--\ref{fig:fmm-g7}.}
\label{fig:fmm-overview}
\end{figure}

%  Each group is broken out in its own detailed figure (fmm_g1 .. fmm_g7).
% ------------------------------------------------------------------------
% ------------------------------------------------------------------------
%  Shared styling for the FMM implementation figures:
%    fmm_overview.tex   -- the whole pipeline, group headlines only
%    fmm_g1..g8.tex     -- one detailed figure per group
%  Idempotent: every figure % ------------------------------------------------------------------------
%  Shared styling for the FMM implementation figures:
%    fmm_overview.tex   -- the whole pipeline, group headlines only
%    fmm_g1..g8.tex     -- one detailed figure per group
%  Idempotent: every figure % ------------------------------------------------------------------------
%  Shared styling for the FMM implementation figures:
%    fmm_overview.tex   -- the whole pipeline, group headlines only
%    fmm_g1..g8.tex     -- one detailed figure per group
%  Idempotent: every figure \input{figures/fmm_style} before its tikzpicture.
%  Needs tikz with the libraries positioning, fit and arrows.meta.
% ------------------------------------------------------------------------

% Entry macros. The mapping runs on from the kernel name in the same paragraph
% rather than on its own line -- a forced break per entry doubles the height.
\providecommand{\fmmH}[1]{{\footnotesize\bfseries #1}\par\vspace{0.5ex}}
\providecommand{\fmmK}[2]{{\ttfamily\scriptsize #1}\,{\tiny\itshape\color{black!62}#2}\par\vspace{0.35ex}}
\providecommand{\fmmN}[1]{{\tiny\color{black!62}#1}\par\vspace{0.35ex}}

% Legend, shared by the overview and by any detail figure that wants it.
\providecommand{\fmmlegend}{{\scriptsize
   \tikz[baseline=-0.5ex]\node[draw, line width=0.6pt, fill=white,
        minimum width=4mm, minimum height=2.6mm, inner sep=0pt]{};\;GPU kernel
   \quad
   \tikz[baseline=-0.5ex]\node[draw=black!55, line width=0.5pt,
        dash pattern=on 1.4pt off 1.2pt, fill=black!7,
        minimum width=4mm, minimum height=2.6mm, inner sep=0pt]{};\;host (CPU) step
   \quad {\itshape italics} $=$ parallel mapping
   \quad \texttt{cub::} $=$ CUB device-wide primitive (GPU)}}

\tikzset{
  fmmgpu/.style  ={draw, line width=0.6pt, fill=white, align=left, inner sep=1.5mm},
  fmmcpu/.style  ={fmmgpu, draw=black!55, dash pattern=on 1.4pt off 1.2pt,
                   fill=black!7},
  fmmnote/.style ={fmmgpu, draw=black!45, fill=black!4, rounded corners=2pt},
  fmmtag/.style  ={font=\tiny\sffamily, text=black!55, fill=white,
                   inner sep=0.5mm, anchor=north east},
  fmmflow/.style ={->, line width=0.7pt, black!70},
  fmmgrp/.style  ={draw=black!45, line width=0.9pt, rounded corners=2.5pt,
                   inner sep=2.4mm},
  fmmlbl/.style  ={font=\small, fill=white, inner sep=1mm, anchor=west},
  fmmedge/.style ={font=\tiny\sffamily, text=black!60, inner sep=0.7mm,
                   fill=white},
  fmmband/.style ={font=\scriptsize\itshape, text=black!60, inner sep=1mm,
                   fill=white},
}
 before its tikzpicture.
%  Needs tikz with the libraries positioning, fit and arrows.meta.
% ------------------------------------------------------------------------

% Entry macros. The mapping runs on from the kernel name in the same paragraph
% rather than on its own line -- a forced break per entry doubles the height.
\providecommand{\fmmH}[1]{{\footnotesize\bfseries #1}\par\vspace{0.5ex}}
\providecommand{\fmmK}[2]{{\ttfamily\scriptsize #1}\,{\tiny\itshape\color{black!62}#2}\par\vspace{0.35ex}}
\providecommand{\fmmN}[1]{{\tiny\color{black!62}#1}\par\vspace{0.35ex}}

% Legend, shared by the overview and by any detail figure that wants it.
\providecommand{\fmmlegend}{{\scriptsize
   \tikz[baseline=-0.5ex]\node[draw, line width=0.6pt, fill=white,
        minimum width=4mm, minimum height=2.6mm, inner sep=0pt]{};\;GPU kernel
   \quad
   \tikz[baseline=-0.5ex]\node[draw=black!55, line width=0.5pt,
        dash pattern=on 1.4pt off 1.2pt, fill=black!7,
        minimum width=4mm, minimum height=2.6mm, inner sep=0pt]{};\;host (CPU) step
   \quad {\itshape italics} $=$ parallel mapping
   \quad \texttt{cub::} $=$ CUB device-wide primitive (GPU)}}

\tikzset{
  fmmgpu/.style  ={draw, line width=0.6pt, fill=white, align=left, inner sep=1.5mm},
  fmmcpu/.style  ={fmmgpu, draw=black!55, dash pattern=on 1.4pt off 1.2pt,
                   fill=black!7},
  fmmnote/.style ={fmmgpu, draw=black!45, fill=black!4, rounded corners=2pt},
  fmmtag/.style  ={font=\tiny\sffamily, text=black!55, fill=white,
                   inner sep=0.5mm, anchor=north east},
  fmmflow/.style ={->, line width=0.7pt, black!70},
  fmmgrp/.style  ={draw=black!45, line width=0.9pt, rounded corners=2.5pt,
                   inner sep=2.4mm},
  fmmlbl/.style  ={font=\small, fill=white, inner sep=1mm, anchor=west},
  fmmedge/.style ={font=\tiny\sffamily, text=black!60, inner sep=0.7mm,
                   fill=white},
  fmmband/.style ={font=\scriptsize\itshape, text=black!60, inner sep=1mm,
                   fill=white},
}
 before its tikzpicture.
%  Needs tikz with the libraries positioning, fit and arrows.meta.
% ------------------------------------------------------------------------

% Entry macros. The mapping runs on from the kernel name in the same paragraph
% rather than on its own line -- a forced break per entry doubles the height.
\providecommand{\fmmH}[1]{{\footnotesize\bfseries #1}\par\vspace{0.5ex}}
\providecommand{\fmmK}[2]{{\ttfamily\scriptsize #1}\,{\tiny\itshape\color{black!62}#2}\par\vspace{0.35ex}}
\providecommand{\fmmN}[1]{{\tiny\color{black!62}#1}\par\vspace{0.35ex}}

% Legend, shared by the overview and by any detail figure that wants it.
\providecommand{\fmmlegend}{{\scriptsize
   \tikz[baseline=-0.5ex]\node[draw, line width=0.6pt, fill=white,
        minimum width=4mm, minimum height=2.6mm, inner sep=0pt]{};\;GPU kernel
   \quad
   \tikz[baseline=-0.5ex]\node[draw=black!55, line width=0.5pt,
        dash pattern=on 1.4pt off 1.2pt, fill=black!7,
        minimum width=4mm, minimum height=2.6mm, inner sep=0pt]{};\;host (CPU) step
   \quad {\itshape italics} $=$ parallel mapping
   \quad \texttt{cub::} $=$ CUB device-wide primitive (GPU)}}

\tikzset{
  fmmgpu/.style  ={draw, line width=0.6pt, fill=white, align=left, inner sep=1.5mm},
  fmmcpu/.style  ={fmmgpu, draw=black!55, dash pattern=on 1.4pt off 1.2pt,
                   fill=black!7},
  fmmnote/.style ={fmmgpu, draw=black!45, fill=black!4, rounded corners=2pt},
  fmmtag/.style  ={font=\tiny\sffamily, text=black!55, fill=white,
                   inner sep=0.5mm, anchor=north east},
  fmmflow/.style ={->, line width=0.7pt, black!70},
  fmmgrp/.style  ={draw=black!45, line width=0.9pt, rounded corners=2.5pt,
                   inner sep=2.4mm},
  fmmlbl/.style  ={font=\small, fill=white, inner sep=1mm, anchor=west},
  fmmedge/.style ={font=\tiny\sffamily, text=black!60, inner sep=0.7mm,
                   fill=white},
  fmmband/.style ={font=\scriptsize\itshape, text=black!60, inner sep=1mm,
                   fill=white},
}


\begin{figure}[p]
\centering
\begin{tikzpicture}[
    >={Latex[length=1.9mm]},
    b/.style ={fmmgpu, text width=4.4cm, align=center, inner sep=1.4mm,
               font=\small},
    bc/.style={fmmcpu, text width=4.4cm, align=center, inner sep=1.4mm,
               font=\small},
    w/.style ={b,  text width=9.7cm},
    wc/.style={bc, text width=9.7cm},
]

% ---------- legend ----------
\node[anchor=north, align=center] at (4.99,0.5) {{\scriptsize
   \tikz[baseline=-0.5ex]\node[draw, line width=0.6pt, fill=white,
        minimum width=4mm, minimum height=2.6mm, inner sep=0pt]{};\;GPU kernel
   \qquad
   \tikz[baseline=-0.5ex]\node[draw=black!55, line width=0.5pt,
        dash pattern=on 1.4pt off 1.2pt, fill=black!7,
        minimum width=4mm, minimum height=2.6mm, inner sep=0pt]{};\;host (CPU) step}};

% =============== 1  source tree construction ===============
\node[wc, anchor=north west] (a1) at (0,-1.0) {Setup};
\node[w,  anchor=north west] (a2) at ([yshift=-1.8mm]a1.south west) {LBVH build};
\node[w,  anchor=north west] (a3) at ([yshift=-1.8mm]a2.south west) {Box nodes and box leaves};
\node[wc, anchor=north west] (a4) at ([yshift=-1.8mm]a3.south west) {Moment storage};
\node[fmmgrp, fit=(a1)(a2)(a3)(a4)] (G1) {};

% =============== 2  upward pass ===============
% Every group is anchored back to x = 0, never to the previous PANEL edge --
% otherwise each group drifts 2.4 mm (one fmmgrp inner sep) further left.
\coordinate (p2) at (0,0 |- G1.south);
\node[b, anchor=north west] (b1) at ([yshift=-8mm]p2) {P2M};
\node[b, anchor=north west] (b2) at ([xshift=6mm]b1.north east) {M2M};
\node[fmmgrp, fit=(b1)(b2)] (G2) {};

% =============== 3  target tree construction ===============
\coordinate (p3) at (0,0 |- G2.south);
\node[w,  anchor=north west] (c1) at ([yshift=-23mm]p3) {LBVH build, box nodes and box leaves};
\node[w,  anchor=north west] (c2) at ([yshift=-1.8mm]c1.south west) {L2L depth worklist};
\node[wc, anchor=north west] (c3) at ([yshift=-1.8mm]c2.south west) {Local-expansion storage};
\node[fmmgrp, fit=(c1)(c2)(c3)] (G3) {};

% =============== 4  dual-tree traversal ===============
\coordinate (p4) at (0,0 |- G3.south);
\node[bc, anchor=north west] (d1) at ([yshift=-8mm]p4) {Host driver};
\node[b,  anchor=north west] (d2) at ([xshift=6mm]d1.north east) {\texttt{\small fmm\_pair\_kernel}};
\node[fmmgrp, fit=(d1)(d2)] (G4) {};

% =============== 5 / 6  the two interaction lists ===============
% M2L and L2L share one host function (fmm_translations<P>), so they share a panel.
\coordinate (p5) at (0,0 |- G4.south);
\node[b, anchor=north west] (e1) at ([yshift=-11mm]p5) {Near-list CSR};
\node[b, anchor=north west] (f1) at ([xshift=6mm]e1.north east) {M2L accumulate};
\node[b, anchor=north west] (f2) at ([yshift=-1.8mm]f1.south west) {L2L level sweep};
\node[fmmgrp, fit=(e1)] (G5) {};
\node[fmmgrp, fit=(f1)(f2)] (G6) {};

% =============== 7  fused evaluation ===============
\coordinate (p7) at (0,0 |- G6.south);
\node[w,  anchor=north west] (h1) at ([yshift=-11mm]p7) {Fused L2P $+$ P2P leaf kernel};
\node[wc, anchor=north west] (h2) at ([yshift=-1.8mm]h1.south west) {Host reduction};
\node[fmmgrp, fit=(h1)(h2)] (G7) {};

% ---------- group labels on the panel borders ----------
\node[fmmlbl] at ([xshift=4mm]G1.north west) {\bfseries 1\ \ Source tree construction};
\node[fmmlbl] at ([xshift=4mm]G2.north west) {\bfseries 2\ \ Upward pass};
\node[fmmlbl] at ([xshift=4mm]G3.north west) {\bfseries 3\ \ Target tree construction};
\node[fmmlbl] at ([xshift=4mm]G4.north west) {\bfseries 4\ \ Dual-tree traversal};
\node[fmmlbl] at ([xshift=4mm]G5.north west) {\bfseries 5\ \ P2P CSR};
\node[fmmlbl] at ([xshift=4mm]G6.north west) {\bfseries 6\ \ Downward pass};
\node[fmmlbl] at ([xshift=4mm]G7.north west) {\bfseries 7\ \ Fused evaluation};

% ---------- flow ----------
\foreach \x/\y in {a1/a2, a2/a3, a3/a4, c1/c2, c2/c3, f1/f2, h1/h2}
  \draw[fmmflow] (\x.south) -- (\y.north);
\draw[fmmflow] (G1.south) -- (G2.north);
\draw[fmmflow] (G2.south) -- (G3.north);
\draw[fmmflow] (G3.south) -- (G4.north);

% Connectors run down the right of each panel: the group labels sit on the
% top-left of the borders, so a centered arrow would strike through them.
\coordinate (x5) at ([xshift=-8mm]G5.north east);
\coordinate (x6) at ([xshift=-8mm]G6.north east);

% the single traversal pass forks into the two interaction lists
\draw[fmmflow] (x5 |- G4.south) -- node[fmmedge, pos=0.55] {P2P pair list} (x5);
\draw[fmmflow] (x6 |- G4.south) -- node[fmmedge, pos=0.55] {M2L pair list} (x6);

% and both rejoin in the fused leaf kernel
\draw[fmmflow] (x5 |- G5.south) -- node[fmmedge, pos=0.5] {near field} (x5 |- G7.north);
\draw[fmmflow] (x6 |- G6.south) -- node[fmmedge, pos=0.5] {far field} (x6 |- G7.north);

% ---------- lifecycle band (drawn last so it masks the arrow behind it) ------
\draw[black!45, line width=0.5pt, dash pattern=on 1.6pt off 1.6pt]
     ([xshift=-6mm,yshift=9.5mm]G3.north west) -- ([xshift=6mm,yshift=9.5mm]G3.north east);
\node[fmmband, anchor=south west] at ([xshift=-4mm,yshift=10.1mm]G3.north west)
     {once per solve};
\node[fmmband, anchor=north west] at ([xshift=-4mm,yshift=8.9mm]G3.north west)
     {once \emph{per energy term}: $E_{\text{pol}}$, then $E_{\text{ion}}$};

\end{tikzpicture}
\caption[FMM implementation: pipeline overview]{%
    \textbf{Overview of the FMM implementation.} Solid boxes are GPU kernels,
    dashed boxes host steps. Each group is broken out at the level of individual
    kernels in Figures~\ref{fig:fmm-g1}--\ref{fig:fmm-g7}.}
\label{fig:fmm-overview}
\end{figure}
